
\subsubsection{2.1 Dataset Documentation Frameworks}

Gebru et al. (2021) proposed Datasheets for Datasets, a template covering seven documentation sections (Motivation, Composition, Collection Process, Preprocessing/Cleaning/Labeling, Uses, Distribution, Maintenance). Bender and Friedman (2018) introduced Data Statements for NLP, a domain-specific standard addressing speaker demographics and curation rationale. Pushkarna et al. (2022) developed Data Cards, a structured documentation format validated across deployments at Google and operationally related to HuggingFace's dataset card format.

The Data Transparency Scorecard \citep{Rondina2025DTS} converts the Datasheets taxonomy into a six-section rubric with inverse-frequency weights derived from empirical section rarity across 100 manually scored datasets. We adopt the DTS framework directly in this work. It should be noted that Rondina et al. (2025) has not been independently confirmed in Semantic Scholar or other literature databases at the time of writing; all DTS weight values used in this paper are dependent on the accuracy of that source.

These frameworks specify what information should be documented but do not themselves measure compliance at scale across repositories.

\subsubsection{2.2 Empirical Studies of Documentation Quality}

Paullada et al. (2021) surveyed documentation limitations in ML dataset development, identifying widespread gaps and a research culture that undervalues documentation investment. Sambasivan et al. (2021) reported that 92\% of surveyed AI practitioners experience compounding negative effects from poor data quality -- termed ''Data Cascades'' -- driven in part by documentation underinvestment. Koch et al. (2021) quantified dataset usage concentration across ML communities from 2015 to 2020.

Rondina et al. (2025) is the most directly related predecessor: manual DTS scoring of 100 datasets across four repositories, reporting near-zero coverage for Preprocessing and Collection sections and higher coverage for Motivation. Our work extends this to automated scoring at 7.58 times the sample size, though covering three repositories rather than four.

\subsubsection{2.3 Automated Metadata Analysis}

Oreamuno et al. (2024) performed binary field presence checks on 6,758 HuggingFace Hub dataset cards, reporting that 71.52\% have substantial undocumented sections. Their analysis is single-repository and unweighted, treating all fields equally regardless of DTS importance. Lhoest et al. (2021) introduced the HuggingFace \texttt{datasets} library with the \texttt{card\textbackslash \_data} structured YAML API, which defines the upper bound on what automated DTS scoring from HuggingFace can retrieve.

No prior work applies DTS inverse-frequency weighting in an automated cross-repository pipeline or measures per-section API coverage against DTS importance weights.
